% =============================================================================
% Section 10: Conclusion
% =============================================================================
\section{Conclusion}
\label{sec:conclusion}

% ---------------------------------------------------------------------------
% 10.1 Summary of Contributions
% ---------------------------------------------------------------------------
\subsection{Summary of Contributions}
\label{sec:conc:summary}

This thesis has presented a three-stage complexity ladder for meningioma
growth forecasting from longitudinal multi-modal 3D MRI. Rather than
proposing a single complex pipeline and reporting its metrics, we
systematically evaluate whether each successive level of model
sophistication provides measurable out-of-sample improvement at $N = 31$--$58$
patients. The key contributions and findings are:

\begin{enumerate}[leftmargin=*]
  \item \textbf{The first out-of-sample meningioma growth prediction
    benchmark.} Using LOPO-CV on a longitudinal cohort, we provide the
    first reported $R^2$, MAE, and calibration metrics for meningioma
    volume prediction. The volumetric baseline (Stage~1: ScalarGP, LME,
    HGP) establishes a reference point that, to our knowledge, did not
    previously exist in the literature. The HGP with population linear mean
    from the LME provides calibrated uncertainty with automatic shrinkage
    toward the population trend for patients with sparse observations.

  \item \textbf{A latent severity framework for growth characterisation.}
    Stage~2 formulates meningioma growth heterogeneity as a latent variable
    problem connected to Item Response Theory and the reduced Gompertz model
    \cite{vaghi2020reduced}. A single patient-specific severity parameter
    $s_i \in [0, 1]$ captures inter-patient variability while respecting
    the sample size constraint of 2--3 parameters per predictor. The
    framework provides interpretable patient stratification aligned with
    the three growth subgroups identified by Engelhardt et al.\
    \cite{engelhardt2023gompertz}.

  \item \textbf{Systematic assessment of representation learning value.}
    By placing the full BrainSegFounder + LoRA + SDP pipeline as Stage~3,
    we directly quantify whether 768-dimensional deep features (compressed
    via PCA to 3--6 dimensions) capture growth-relevant information beyond
    what volume and severity already explain. The marginal benefit of LoRA
    adaptation for feature quality ($\Delta R^2 = +0.016$ on semantic
    probes) motivates scepticism that the full representation learning
    pipeline is necessary for this task.

  \item \textbf{Variance decomposition as an analytical framework.}
    The cross-cutting analysis decomposes LOPO-CV prediction variance across
    the three stages using paired permutation tests and BCa bootstrap CIs.
    This provides not only a model selection criterion but also a
    methodological template for evaluating complexity--benefit tradeoffs in
    small-sample longitudinal prediction problems more broadly.

  \item \textbf{Empirically-driven methodology revision.} Rigorous
    evaluation revealed that shape disentanglement fails ($R^2 \leq 0.11$),
    meningioma location is temporally static, and sub-compartment volume
    targets are label-dependent and clinically irrelevant. These findings
    motivated a principled simplification and the reframing from a linear
    pipeline to the complexity ladder.
\end{enumerate}

% ---------------------------------------------------------------------------
% 10.2 Limitations
% ---------------------------------------------------------------------------
\subsection{Limitations}
\label{sec:conc:limitations}

\subsubsection{Cohort Size}

The longitudinal cohort ($N = 31$--$58$ patients, $\sim$112~observations)
is the primary bottleneck. With approximately 57.6\% of patients having
only 2~studies, per-patient growth trajectories are sparse. The GP and LME
frameworks handle this via automatic shrinkage toward the population mean,
but predictions for 2-observation patients are necessarily dominated by
population-level dynamics rather than individual-specific growth rates.
The sample size also limits the statistical power of the variance
decomposition: Varoquaux \cite{varoquaux2018cv} showed that CV error
estimates have standard deviations of $\sim$$\pm$10\% at $N = 100$, and
uncertainty is larger still at $N = 33$.

\subsubsection{Ordinal Time}

At the time of writing, exact acquisition timestamps for the longitudinal
cohort are pending. If only ordinal timepoint indices are available (scan~1,
scan~2, scan~3, \ldots) rather than calendar dates, the temporal models
operate on a surrogate time axis. This does not invalidate the complexity
ladder comparison (all models use the same time representation) but limits
the clinical interpretability of growth rate estimates and length-scale
hyperparameters.

\subsubsection{LoRA Marginality and Stage~3 Viability}

The marginal benefit of LoRA adaptation ($\Delta R^2 = +0.016$ for
semantic probing) suggests that the entire Stage~3 pipeline may provide
limited value at the current sample size. This is, however, a finding
rather than a failure: establishing that simple volumetric models suffice
is a valuable negative result that can save computational resources in
clinical deployment.

\subsubsection{Single-Site Longitudinal Data}

The longitudinal cohort is collected from hospitals in a single region
(Andalusia, Spain), which may limit generalisability to other populations,
scanner manufacturers, and acquisition protocols. ComBat harmonisation
addresses scanner batch effects but cannot correct for biological population
differences.

\subsubsection{Absence of Clinical Validation}

The pipeline has not been validated against clinical outcomes (intervention
decisions, actual tumour behaviour at longer follow-up). Volume prediction
$R^2$ and calibration metrics are necessary but not sufficient for clinical
utility. A positive growth prediction does not, by itself, determine the
appropriate clinical action.

% ---------------------------------------------------------------------------
% 10.3 Future Work
% ---------------------------------------------------------------------------
\subsection{Future Work}
\label{sec:conc:future}

\subsubsection{Cohort Expansion to $N \geq 50$}

The most impactful improvement would be expanding the longitudinal cohort.
Multi-centre collaborations targeting $N \geq 50$ patients would: (i)~tighten
the variance decomposition CIs, enabling confident $\Delta R^2$ conclusions;
(ii)~potentially unlock value in Stage~2 and~3 by providing sufficient data
for more complex models; (iii)~enable subgroup analyses by tumour grade,
location, and patient demographics.

\subsubsection{Exact Timestamps}

Obtaining exact acquisition dates would enable physically meaningful
growth rate estimation, improve GP length-scale interpretability, and allow
direct comparison with the volumetric growth rates reported in the clinical
literature \cite{hashimoto2012natural,nakamura2003natural,
engelhardt2023gompertz}.

\subsubsection{Clinical Decision Support Integration}

The GP posterior's calibrated uncertainty is directly suitable for clinical
decision support: the probability that a tumour will exceed a volume
threshold $V_{\text{critical}}$ within a specified time horizon can be
computed analytically from the Gaussian predictive distribution:
\begin{equation}
  P(V(t_{\text{future}}) > V_{\text{critical}}) =
    1 - \Phi\!\left(\frac{\log(V_{\text{critical}} + 1) - \mu_*(t_{\text{future}})}
    {\sigma_*(t_{\text{future}})}\right),
  \label{eq:clinical_decision}
\end{equation}
where $\Phi$ is the standard normal CDF and $(\mu_*, \sigma_*)$ are the GP
posterior mean and standard deviation in log-volume space. Integrating this
into a clinical workflow would require prospective validation and regulatory
approval.

\subsubsection{Severity-Guided Surveillance}

If the severity model (Stage~2) provides robust patient stratification,
it could inform surveillance scheduling: high-severity patients receive
shorter follow-up intervals, while low-severity patients can be monitored
less frequently. This requires validation of the severity--growth
relationship at longer follow-up horizons.

\subsubsection{Foundation Model Scaling}

BrainSegFounder is a relatively small foundation model (BSF-Tiny, 62M
parameters). Larger variants or more recent foundation models pretrained on
broader datasets may provide stronger initial features, potentially shifting
the Stage~3 $\Delta R^2$ from marginal to significant. However, this
hypothesis must be tested empirically rather than assumed.
